% Figure: energy-flow (Sankey-style) diagram for a refrigerator. Qc drawn from
% cold space + W input combine into Qh rejected at the condenser.
\begin{figure}[htbp]
\centering
\begin{tikzpicture}[font=\small]
  % cold space heat in (Qc ~ 50 W): band height 1.0
  \fill[engblue!35] (0,3.0) rectangle (3.0,4.0);
  \node[font=\scriptsize, anchor=east] at (-0.1,3.5) {$Q_c \approx 50\,\text{W}$};
  \node[engblue, font=\scriptsize] at (1.5,3.5) {from cold space};
  % work input (W ~ 18 W): band height 0.36
  \fill[englime!45] (0,2.4) rectangle (3.0,2.76);
  \node[font=\scriptsize, anchor=east] at (-0.1,2.58) {$W \approx 18\,\text{W}$};
  \node[font=\scriptsize] at (1.5,2.58) {compressor};
  % merge into Qh (height 1.36) leaving to the right at condenser
  \fill[engred!35] (3.0,2.4) -- (5.2,2.4) -- (5.2,4.0) -- (3.0,4.0) -- cycle;
  \node[engred, font=\scriptsize] at (4.1,3.2) {$Q_h \approx 68\,\text{W}$};
  \node[engred, font=\scriptsize, anchor=west] at (5.3,3.2) {rejected at condenser};
  % boundary marker
  \draw[engblack, very thick, dashed] (3.0,1.8) -- (3.0,4.4);
  \node[enggray, font=\scriptsize, anchor=north] at (3.0,1.8) {refrigerator};
\end{tikzpicture}
\caption[Energy flow through a household refrigerator]{Energy flows through a
refrigerator at steady state. The unit pumps heat $Q_c$ out of the cold cabinet
and the compressor adds work $W$; the first law forces the condenser to reject
the sum $Q_h = Q_c + W$ to the kitchen. A refrigerator therefore warms the room
it stands in, by more than the electrical power it draws. The widths are drawn
to scale for the worked estimate of roughly $50\,\text{W}$ pumped, $18\,\text{W}$
of electrical input, and $68\,\text{W}$ rejected; the coefficient of
performance $Q_c/W \approx 2.8$ is the ratio of the two left-hand bands.}
\label{fig:vol04ch01:refrigerator-sankey}
\end{figure}
